\documentclass[11pt]{article}

\usepackage[T1]{fontenc}
\usepackage{lmodern}
\usepackage[margin=0.93in]{geometry}
\usepackage{amsmath,amssymb,amsthm,mathtools}
\usepackage{microtype}
\usepackage{graphicx}
\usepackage[hidelinks]{hyperref}
\usepackage{enumitem}

\setlist{nosep}
\newtheorem{theorem}{Theorem}
\newtheorem{lemma}{Lemma}
\newtheorem{proposition}{Proposition}
\theoremstyle{remark}
\newtheorem{remark}{Remark}

\newcommand{\A}{\mathcal A}
\newcommand{\C}{C(X)}
\newcommand{\Fix}{\operatorname{Fix}}
\newcommand{\dist}{\operatorname{dist}}

\title{Diliberto--Straus Convergence Without Closedness\\for Two-Point Fibres}
\author{Independent research note}
\date{August 17, 2026}

\begin{document}
\maketitle

\begin{abstract}
Asgarova and Ismailov ask whether their Diliberto--Straus residual algorithm
converges to the uniform approximation error without assuming that the sum
of the two algebras is closed.  We prove a positive answer whenever every
quotient fibre of both algebras has at most two points.  Under the paper's
$C$-property, these fibres are the orbits of continuous involutions, so the
midrange maps become linear.  A Laurent-polynomial calculation then gives
an $O(n^{-1/2})$ decay estimate on each summand.  We also construct two such
algebras whose sum is not closed, showing that the result genuinely passes
beyond the source theorem.
\end{abstract}

\section*{Packet status}

\begin{description}[leftmargin=3.0cm,style=nextline]
\item[Run] \texttt{fa\_banach\_001}, lane 13.
\item[Source] arXiv:1603.07073, Theorem~2.3 and the question on PDF page~7 \cite{AI}.
\item[Classification] Candidate strong partial result, likely valid subject to human review.
\item[Scope] A complete answer for the two-point-fibre class; no claim for arbitrary fibres.
\end{description}

\begin{center}
\includegraphics[width=0.94\textwidth]{figures/source_question.png}
\end{center}

\section{Set-up and statement}

Let $X$ be compact Hausdorff and let $\A_1,\A_2$ be closed unital real
subalgebras of $\C$.  Write $s_i:X\to X_i$ for the quotient maps induced
by equality of all functions in $\A_i$.  The source assumes the
\emph{$C$-property}: for every $h\in\C$, the maximum and minimum of $h$ on
each $s_i$-fibre vary continuously on $X_i$.  Its proximity maps are the
fibre midranges
\[
 F_i h(x)=\frac12\left(
 \max_{s_i(y)=s_i(x)}h(y)+\min_{s_i(y)=s_i(x)}h(y)\right).
\]
Starting with $h_1=h$, the algorithm alternately sets
\[
 h_{2n}=h_{2n-1}-F_1h_{2n-1},\qquad
 h_{2n+1}=h_{2n}-F_2h_{2n}.
\]
Put $E(h)=\dist(h,\A_1+\A_2)$.

\begin{theorem}[two-point-fibre convergence]\label{thm:main}
Assume the $C$-property and
\[
 \#s_i^{-1}(s_i(x))\le2\qquad(x\in X,\ i=1,2).
\]
Then, whether or not $\A_1+\A_2$ is closed,
\[
 \lim_{j\to\infty}\|h_j\|_\infty=E(h)
 \qquad(h\in\C).
\]
More precisely, there are continuous involutions $\sigma,\tau:X\to X$
such that $\A_1=\Fix S_\sigma$, $\A_2=\Fix S_\tau$, where
$S_\sigma h=h\circ\sigma$.  If
\[
 T=\frac{I-S_\tau}{2}\frac{I-S_\sigma}{2},
\]
then $h_{2n+1}=T^nh$, and for every $v\in\A_2$,
\begin{equation}\label{eq:rate}
 \|T^nv\|_\infty\le \frac{\|v\|_\infty}{\sqrt n}
 \qquad(n\ge1).
\end{equation}
\end{theorem}

\section*{Idea of the proof}

There are two rigidities hidden in the two-point hypothesis.  First, the
source's $C$-property forces the operation that swaps the two members of a
fibre to be continuous.  Each midrange map is therefore the average of the
identity with an isometric involution.  Second, alternating the complementary
half-differences produces a fixed linear contraction $T$.  On a function
fixed by the second involution, every iterate of $T$ is a Laurent polynomial
in the product of the two involutions.  Its coefficients are discrete
derivatives of a symmetric binomial distribution and have total variation
$O(n^{-1/2})$.  The limiting residual norm is consequently a seminorm that
vanishes on both algebras and their closure.  Since every residual is in the
same quotient class as the initial function, this seminorm must equal the
distance to the algebraic sum.

\section{Two-point fibres are involution orbits}

\begin{lemma}\label{lem:involution}
Let $s:X\to Y$ be one of the algebra quotients, assume its fibres have at
most two points, and assume the $C$-property.  The map $\sigma$ that swaps
the two points of every doubleton fibre and fixes singleton fibres is a
continuous involution.  The associated algebra is $\Fix S_\sigma$, and its
midrange map is $(I+S_\sigma)/2$.
\end{lemma}

\begin{proof}
For $f\in\C$, let $Mf$ be its fibre midrange, regarded as a continuous
function on $X$ via $s$.  On a singleton or doubleton fibre,
\[
 f(\sigma(x))=2Mf(x)-f(x).
\]
Thus $f\circ\sigma$ is continuous for every $f\in\C$.  Continuous real
functions determine the topology of a compact Hausdorff space, so $\sigma$
is continuous.  It is visibly an involution and its orbits are exactly the
fibres.  By Stone--Weierstrass, the closed unital algebra is all continuous
functions on its Hausdorff quotient; hence it consists exactly of the
$\sigma$-invariant functions.  The displayed identity also gives the
formula for the midrange map.
\end{proof}

This lemma also proves the converse description: a continuous involution
has one- or two-point orbits, its invariant algebra has the $C$-property,
and fibre maximum and minimum are
\[
 \frac{f+f\circ\sigma\pm|f-f\circ\sigma|}{2}.
\]

\section{The decay identity}

Set $U=S_\sigma$, $V=S_\tau$ and
\[
 P=\frac{I-U}{2},\qquad Q=\frac{I-V}{2},\qquad T=QP.
\]
These are contractions on $\C$, $U^2=V^2=I$, and
$h_{2n+1}=T^nh$.  Put
\[
 R=VU,\qquad D=\frac{R-R^{-1}}4,
 \qquad L=\frac{2I+R+R^{-1}}4.
\]
Here $R$ and every $R^k$ are surjective isometries.

If $v\in\Fix V$, a direct expansion gives $Tv=Dv$.  The range of $T$ lies
in the anti-invariant space $\{x:Vx=-x\}$, and on that space another direct
expansion gives $Tx=Lx$.  Therefore
\begin{equation}\label{eq:LD}
 T^nv=L^{n-1}Dv\qquad(n\ge1,\ v\in\Fix V).
\end{equation}
For $m=n-1$, the binomial theorem yields
\[
 L^m=\sum_{k=-m}^{m}p_m(k)R^k,
 \qquad
 p_m(k)=4^{-m}\binom{2m}{m+k},
\]
where $p_m(k)=0$ outside $[-m,m]$.  The coefficient of $R^k$ in $L^mD$ is
\[
 q_m(k)=\frac{p_m(k-1)-p_m(k+1)}4.
\]
The sequence $p_m$ is symmetric and unimodal, so its two tails telescope:
\begin{equation}\label{eq:l1}
 \sum_k|q_m(k)|
 =\frac{p_m(0)+p_m(1)}2
 \le p_m(0)
 =4^{-m}\binom{2m}{m}
 \le\frac1{\sqrt{m+1}}.
\end{equation}
The last elementary central-binomial bound also follows immediately by
induction.  Equations~\eqref{eq:LD}--\eqref{eq:l1}, together with
$\|R^k\|=1$, prove~\eqref{eq:rate}.

\begin{proof}[Proof of Theorem~\ref{thm:main}]
Lemma~\ref{lem:involution} supplies $U,V$.  Since $T$ is linear and
contractive, the decreasing limit
\[
 \rho(h)=\lim_{n\to\infty}\|T^nh\|_\infty
\]
is a seminorm satisfying $\rho(h)\le\|h\|_\infty$.  If $u\in\A_1$, then
$Tu=0$; if $v\in\A_2$, estimate~\eqref{eq:rate} gives $\rho(v)=0$.
Therefore $\rho$ vanishes on $\A_1+\A_2$ and, being 1-Lipschitz, on its
closure.  It follows that
\[
 \rho(h)=\rho(h-w)\le\|h-w\|_\infty
 \quad(w\in\A_1+\A_2),
\]
so $\rho(h)\le E(h)$.  In the other direction, every update subtracts an
element of one algebra, hence $h-T^nh\in\A_1+\A_2$ and
$E(h)\le\|T^nh\|_\infty$.  Passing to the limit proves $\rho(h)=E(h)$.
The norms of the entire residual sequence are decreasing, while its odd
subsequence has this limit, so the same is true of all $\|h_j\|_\infty$.
\end{proof}

\section{A genuinely nonclosed example}

The closedness hypothesis is not automatic in Theorem~\ref{thm:main}.
Let $N_j\to\infty$ be even integers and let $X$ be the one-point
compactification of the disjoint union of cycles
$C_{N_j}=\mathbb Z/N_j\mathbb Z$.  On every cycle define
\[
 \sigma(k)=-k,\qquad \tau(k)=1-k,
\]
and fix the point at infinity.  These are continuous involutions.  Let
$\A_1=\Fix S_\sigma$ and $\A_2=\Fix S_\tau$.

On $C_N$, with $\theta=2\pi/N$, put
\[
 u_N(k)=\cos(\theta k),\qquad
 v_N(k)=-\cos\bigl(\theta(k-\tfrac12)\bigr).
\]
Then $u_N\in\A_1|_{C_N}$, $v_N\in\A_2|_{C_N}$, and
\begin{equation}\label{eq:witness}
 \|u_N+v_N\|_\infty\le\frac\pi N.
\end{equation}
Define $y(\infty)=0$ and
$y|_{C_{N_j}}=u_{N_j}+v_{N_j}$.  By~\eqref{eq:witness}, $y$ is continuous,
and truncating to finitely many cycles shows
$y\in\overline{\A_1+\A_2}$.

Suppose $y=a+b$ with $a\in\A_1$, $b\in\A_2$.  Shift a constant between
$a$ and $b$ so that both vanish at infinity.  On each cycle,
$a-u_N=-(b-v_N)$ is fixed by both involutions.  The involutions generate
the transitive dihedral action because their product is a one-step
rotation, so their common fixed functions on $C_N$ are constants.  Since
$N$ is even, $u_N$ takes the values $1$ and $-1$ and has distance one from
the constants.  Hence the restriction of $a$ to every cycle has distance
one from the constants, contradicting $\|a|_{C_{N_j}}\|_\infty\to0$.
Thus $y\notin\A_1+\A_2$, and the sum is not closed.  Nevertheless
Theorem~\ref{thm:main} proves convergence for every $h\in C(X)$.

\section{Boundary of the result}

For fibres of size at least three, the midrange depends on the currently
active maximum and minimum and is no longer a fixed average.  The residual
iteration is then nonlinear, so $\rho(h)=\lim\|h_{2n+1}\|$ need not be a
seminorm and the Laurent cancellation above has no fixed isometry on which
to act.  Eight focused routes were audited: bounded larger fibres, the
source's open-bolt correction, later infinite-bolt duality \cite{AI2}, an
abstract asymptotic functional, and long finite gadgets.  The first four
all encounter this switching-extrema obstruction, while long two-point
gadgets are ruled out by~\eqref{eq:l1}.  No full proof or counterexample for
arbitrary fibres is claimed.

The exact coefficient identity was checked for $0\le m\le100$, and all
operator formulas, the rate, and the nonclosed witnesses were independently
tested on finite dihedral cycles.  A bounded literature audit on August 17,
2026 found no later explicit resolution of the source question or the
two-point-fibre theorem.  The result should be treated as a candidate new
partial answer pending expert review.

\begin{thebibliography}{9}
\bibitem{AI}
A.~Kh.~Asgarova and V.~E.~Ismailov,
\emph{On the Diliberto--Straus algorithm for the uniform approximation by a
sum of two algebras}, arXiv:1603.07073 (2016).

\bibitem{AI2}
A.~Kh.~Asgarova and V.~E.~Ismailov,
\emph{A note on approximation by a sum of two algebras},
arXiv:2405.09815 (2024).
\end{thebibliography}

\end{document}
